\clearpage
\phantomsection
\addcontentsline{toc}{section}{Увод}

\section*{Увод}

\hspace*{2.5em}Теорията на обобщените мрежи (ОМ) предоставя мощен математически апарат за моделиране и симулация на сложни процеси, характеризиращи се с паралелизъм и логическа обвързаност. От самото си зараждане, апаратът на ОМ се утвърждава като гъвкав инструмент в области като индустриалното производство, компютърните системи и изкуствения интелект. С нарастването на сложността на моделите обаче, необходимостта от достъпни и високоефективни софтуерни инструменти за тяхното описание и изследване става все по-осезаема.

Настоящият дисертационен труд е насочен към решаването на ключови проблеми в областта на практическото използване на ОМ. Въпреки съществуващите софтуерни решения, до момента се наблюдава липса на инструменти, които да позволяват работа в онлайн среда, автоматизирано генериране на графични представяния и лесна интеграция в научни публикации. Това ограничава споделянето на модели и затруднява изследователския процес, особено в условията на съвременната колаборация през облачни платформи.

Основната цел на разработката е създаването на софтуерната система \textbf{\textit{OnlineGN}} -- уеб базиран симулатор, който обединява теоретичните основи на ОМ със съвременните уеб технологии. В рамките на труда се предлагат нови алгоритми за автоматична визуализация и експорт на модели, които превръщат \textbf{\textit{OnlineGN}} в комплексна среда за моделиране. Чрез поредица от симулации на класически задачи и реални производствени сценарии в дисертацията се доказва приложимостта и ефективността на разработения инструмент.